%% Relatório do projeto da disciplina CIN0114 (CIn/UFPE), 2026.2.
%% Adaptado do template da disciplina (IEEEtran, conference, bare_conf.tex V1.4b,
%% de Michael Shell, LPPL 1.3). Acrescentados ao preâmbulo do template: babel
%% (hifenização e rótulos em português), mathptmx e fontspec (Times no XeTeX),
%% adjustbox e tikz.
%% As tabelas numéricas entram por \input dos arquivos gerados por
%% scripts/make_report_assets.py; não edite os números aqui.
%% Compilação: XeLaTeX ou tectonic (tectonic relatorio.tex, dentro de report/).
%% No Overleaf, escolher XeLaTeX no menu do compilador.

\documentclass[conference]{IEEEtran}

\usepackage[brazilian]{babel}
% Times, a fonte do IEEEtran: o compilador XeTeX não a aplica sem o fontspec.
\usepackage{mathptmx}
\usepackage{fontspec}
\setmainfont{texgyretermes}[Extension=.otf, UprightFont=*-regular, BoldFont=*-bold, ItalicFont=*-italic, BoldItalicFont=*-bolditalic]
\setmonofont{texgyrecursor}[Extension=.otf, UprightFont=*-regular, BoldFont=*-bold, ItalicFont=*-italic, BoldItalicFont=*-bolditalic, Scale=0.9, Ligatures=NoCommon]
\usepackage{graphicx}
\usepackage{amsmath}
\usepackage{adjustbox}
\usepackage{tikz}
\usetikzlibrary{arrows.meta,positioning,fit}

\addto\captionsbrazilian{%
  \renewcommand{\refname}{Referências}%
  \renewcommand{\figurename}{Fig.}%
  \renewcommand{\tablename}{Tabela}}

\hyphenation{Ma-li-cious Be-nign Non-DoH Packet-Length-Mode Tree-Ex-plai-ner}

% Tabela gerada, reduzida só se passar da largura disponível.
\newcommand{\tabelacoluna}[1]{{\footnotesize\setlength{\tabcolsep}{3pt}%
  \begin{adjustbox}{max width=\columnwidth}\input{tables/#1}\end{adjustbox}}}
% Sublinhado em nome de arquivo, com quebra de linha permitida depois dele.
\newcommand{\us}{\_\allowbreak}
\newcommand{\tabelapagina}[1]{{\footnotesize\setlength{\tabcolsep}{3pt}%
  \begin{adjustbox}{max width=\textwidth}\input{tables/#1}\end{adjustbox}}}

% Quebra de linha no bloco de autores do IEEEtran: quatro autores em duas linhas.
\makeatletter
\newcommand{\linebreakand}{%
  \end{@IEEEauthorhalign}\hfill\mbox{}\par\mbox{}\hfill
  \begin{@IEEEauthorhalign}}
\makeatother

\begin{document}

\title{CIN0114 Técnicas de Ataque e Detecção de Intrusão\\
Reprodução e avaliação de um IDS explicável para ataques DNS over HTTPS}

\author{\IEEEauthorblockN{Amanda Arruda}
\IEEEauthorblockA{\textit{Centro de Informática - UFPE} \\
Recife, Brasil \\
aams2@cin.ufpe.br}
\and
\IEEEauthorblockN{Breno Silva Xavier de Souza}
\IEEEauthorblockA{\textit{Centro de Informática - UFPE} \\
Recife, Brasil \\
bsxs@cin.ufpe.br}
\linebreakand
\IEEEauthorblockN{João Henrique Portela}
\IEEEauthorblockA{\textit{Centro de Informática - UFPE} \\
Recife, Brasil \\
jhpbs@cin.ufpe.br}
\and
\IEEEauthorblockN{Antonio Gonzaga}
\IEEEauthorblockA{\textit{Centro de Informática - UFPE} \\
Recife, Brasil \\
agla@cin.ufpe.br}
}

\maketitle

\begin{abstract}
O protocolo DNS over HTTPS (DoH) cifra as consultas DNS e esconde do administrador da rede os túneis DNS usados para comando e controle e para exfiltração. Os trabalhos anteriores no mesmo conjunto de dados, como o artigo de referência os transcreve, reportam acurácia de até 99,98\%, sem métricas por classe nem explicação da decisão. Zebin, Rezvy e Luo propõem um sistema de detecção de intrusão que classifica fluxos em Non-DoH, Benign-DoH e Malicious-DoH com três Random Forests empilhados e explica as decisões com SHAP. Este trabalho reimplementa o sistema a partir do texto do artigo, reproduz todas as tabelas e figuras de resultado no CIRA-CIC-DoHBrw-2020 (as principais estão neste relatório e as demais no repositório), repete a avaliação em um segundo conjunto de dados com três ferramentas de túnel novas e avalia uma modificação. O artigo admite duas leituras da profundidade das árvores: com profundidade máxima 5, o sistema empilhado não prediz nenhum fluxo Benign-DoH; com profundidade variável, a matriz de confusão do teste fica a 553 fluxos da publicada. Treinado no conjunto original, o sistema detecta 1,81\% dos fluxos das ferramentas novas. Um Random Forest único com peso de classe, sem amostras sintéticas, supera o sistema empilhado em 0,5 ponto percentual de F1 macro nos dois conjuntos.
\end{abstract}

\IEEEpeerreviewmaketitle

\section{Introdução}

O DoH envia a consulta DNS dentro de uma conexão HTTPS a um resolvedor público. A cifra protege a privacidade do usuário e, ao mesmo tempo, tira do firewall e do IDS o nome consultado, que era o que permitia reconhecer um túnel DNS \cite{zebin2022}. O problema investigado é detectar, só com estatísticas de fluxo e sem decifrar o TLS, o tráfego DoH gerado por ferramentas de túnel. Ele interessa porque um host comprometido passa a ter um canal de saída que se confunde com navegação comum.

Segundo o artigo de referência, as soluções anteriores ou dependem de limiares ajustados ao ambiente, ou reportam só acurácia, ou não descrevem o pré-processamento (Seção II de \cite{zebin2022}; Tabela \ref{tab:relacionados}); o artigo se apresenta como uma das primeiras soluções explicáveis para o problema (Seção I). Ele propõe o \textit{Balanced Stacked Random Forest}: três Random Forests treinados em subconjuntos balanceados, combinados por regressão logística, com explicações SHAP \cite{lundberg2017} em um painel interativo.

As contribuições deste trabalho são:
\begin{itemize}
    \item reimplementação do sistema a partir do texto, com as duas leituras da profundidade das árvores que o artigo admite, reportadas lado a lado;
    \item reprodução de todas as tabelas e figuras de resultado do artigo, com a diferença para o valor publicado: as principais neste relatório, as demais no repositório;
    \item avaliação em um segundo conjunto de dados, com ferramentas de túnel ausentes do original;
    \item avaliação com dez seeds, sem os vetores do teste repetidos no treino e com separação por máquina de origem;
    \item uma modificação do sistema e uma avaliação de robustez à fragmentação do fluxo.
\end{itemize}

\section{Trabalhos relacionados}

A Tabela \ref{tab:relacionados} resume os trabalhos que o artigo de referência discute, com a caracterização que ele dá a cada um. Aiello et al. \cite{aiello2013} e Preston \cite{preston2019} tratam do túnel sobre DNS em claro. MontazeriShatoori et al. \cite{montazeri2020} publicam o CIRA-CIC-DoHBrw-2020 e o extrator de atributos DoHLyzer. Banadaki \cite{banadaki2020}, Jafar et al. \cite{jafar2021} e Ahakonye et al. \cite{ahakonye2022} classificam o mesmo conjunto de dados e são os trabalhos com que o artigo compara seus números (Tabela II de \cite{zebin2022}). Ramakrishnan e Rajan \cite{ramakrishnan2022} propõem um IDS com rede neural. Vekshin et al. \cite{vekshin2020} e Behnke et al. \cite{behnke2021} são citados pelo artigo só na introdução e nos trabalhos futuros.

\begin{table}[!t]
\centering
\caption{Trabalhos relacionados, com a contribuição e a limitação que o artigo de referência atribui a cada um.}
\label{tab:relacionados}
\footnotesize\setlength{\tabcolsep}{3pt}
\begin{tabular}{p{1.95cm}p{3.0cm}p{3.0cm}}
\hline
\textbf{Trabalho} & \textbf{Contribuição} & \textbf{Limitação} \\ \hline
Aiello et al. \cite{aiello2013} & PCA e informação mútua para identificar túnel DNS & Limiar depende do ambiente \\
Preston \cite{preston2019} & Aprendizado supervisionado sobre o domínio primário & Não detecta consulta maliciosa no domínio principal \\
Banadaki \cite{banadaki2020} & XGBoost, Gradient Boosting e LightGBM no mesmo conjunto & Pré-processamento não descrito; teste com 4000 amostras \\
Ramakrishnan e Rajan \cite{ramakrishnan2022} & IDS com rede neural e poucos atributos & Cerca de 90\% de acurácia \\
Jafar et al. \cite{jafar2021} & Oito algoritmos clássicos no mesmo conjunto & Só acurácia e tempo, sem métricas por classe \\
Ahakonye et al. \cite{ahakonye2022} & Ensemble com tempo de processamento reduzido & Não apontada pelo artigo \\
\hline
\end{tabular}
\end{table}

\section{Modelo de ameaça}

O artigo não declara um modelo de ameaça. O que segue é inferência da equipe a partir da Fig. 1, da Seção II-A e da Tabela I de \cite{zebin2022}; a Fig. \ref{fig:ameaca} mostra o cenário.

\begin{figure}[!t]
\centering
\begin{tikzpicture}[font=\scriptsize, node distance=5mm and 11mm,
  caixa/.style={draw, rounded corners=1pt, align=center, inner sep=2.5pt, minimum height=8mm},
  seta/.style={-{Latex[length=1.6mm]}, semithick}]
\node[caixa] (host) {Host interno\\comprometido\\(cliente de túnel)};
\node[caixa, right=of host] (res) {Resolvedor\\DoH público};
\node[caixa, right=of res] (aut) {Servidor\\autoritativo\\do atacante};
\draw[seta] (host) -- node[above, inner sep=1.5pt] (https) {HTTPS} (res);
\draw[seta] (res) -- node[above, inner sep=1.5pt] {DNS} (aut);
\node[caixa, below=of https] (ids) {IDS na borda\\(estatísticas de fluxo)};
\draw[dashed, semithick] (ids.north) -- (https.south);
\end{tikzpicture}
\caption{Túnel DNS sobre DoH: os dados saem codificados em consultas DNS, dentro de HTTPS, até o servidor autoritativo que o atacante controla. O IDS observa o fluxo cifrado na borda da rede da organização.}
\label{fig:ameaca}
\end{figure}

\textbf{Ataque.} Canal encoberto para comando e controle e para exfiltração: o cliente de túnel codifica os dados em consultas DNS de um domínio do atacante e as envia por DoH a um resolvedor público, que as encaminha ao servidor autoritativo desse domínio. \textbf{Premissas:} (i) há um host interno comprometido, capaz de executar dns2tcp, DNSCat2 ou iodine em configuração padrão; (ii) o atacante controla o servidor autoritativo do domínio; (iii) a rede permite HTTPS a resolvedores DoH públicos (AdGuard, Cloudflare, Google, Quad9); (iv) o defensor observa passivamente a borda e só tem estatísticas de fluxo; (v) o atacante não conhece o detector e não se adapta a ele, pois o artigo não avalia evasão. A premissa (v) é relaxada na Seção \ref{sec:equipe}.

\section{Sistema proposto pelo artigo de referência}
\label{sec:sistema}

A entrada é um fluxo descrito por 29 atributos numéricos extraídos pelo DoHLyzer (duração, bytes e taxas, e estatísticas de comprimento de pacote, de tempo entre pacotes e de tempo de resposta). A saída é uma de três classes, com a explicação da decisão. A Fig. \ref{fig:pipeline} mostra os componentes.

\begin{figure}[!t]
\centering
\begin{tikzpicture}[font=\scriptsize, node distance=3.5mm and 3mm,
  caixa/.style={draw, rounded corners=1pt, align=center, inner sep=2.5pt, minimum height=6mm},
  seta/.style={-{Latex[length=1.6mm]}, semithick}]
\node[caixa] (dados) {29 atributos\\por fluxo};
\node[caixa, right=of dados] (split) {Split 90/10\\estratificado};
\node[caixa, right=of split] (norm) {Min-max\\(ajuste no treino)};
\node[caixa, below=5mm of norm] (s2) {Subconjunto 2};
\node[caixa, left=of s2] (s1) {Subconjunto 1};
\node[caixa, right=of s2] (s3) {Subconjunto 3};
\node[caixa, below=of s1] (r1) {RF 1};
\node[caixa, below=of s2] (r2) {RF 2};
\node[caixa, below=of s3] (r3) {RF 3};
\node[caixa, below=of r2] (meta) {Regressão logística};
\node[caixa, left=of meta] (shap) {SHAP\\(TreeExplainer)};
\node[caixa, right=of meta] (saida) {Classe};
\draw[seta] (dados) -- (split);
\draw[seta] (split) -- (norm);
\draw[seta] (norm) -- (s2);
\draw[seta] (norm) -- (s1.north east);
\draw[seta] (norm) -- (s3.north);
\draw[seta] (s1) -- (r1);
\draw[seta] (s2) -- (r2);
\draw[seta] (s3) -- (r3);
\draw[seta] (r1) -- (meta.north west);
\draw[seta] (r2) -- (meta);
\draw[seta] (r3) -- (meta.north east);
\draw[seta] (meta) -- (saida);
\draw[seta] (r1) -- (shap);
\end{tikzpicture}
\caption{Sistema do artigo. Cada subconjunto recebe um terço do Non-DoH, todo o Malicious-DoH e o Benign-DoH aumentado com SMOTE até igualar o Malicious-DoH.}
\label{fig:pipeline}
\end{figure}

\textbf{Pré-processamento.} Split de 90\% para treino e 10\% para teste e normalização min-max ajustada no treino (Seção III de \cite{zebin2022}):
\begin{equation}
x' = \frac{x - x_{\min}}{x_{\max} - x_{\min}}.
\label{eq:minmax}
\end{equation}

\textbf{Balanceamento.} O Non-DoH do treino é dividido em três partes; cada subconjunto recebe uma parte, todos os fluxos Malicious-DoH e o Benign-DoH aumentado com SMOTE. O artigo justifica a divisão pelo desbalanceamento das classes e por um tempo de treino três vezes menor, que não mede (Seções I e III-B). Declara a razão 15:12:12 entre as classes de cada subconjunto; a obtida é 14,3:12:12.

\textbf{Classificação.} Um Random Forest por subconjunto, com 10 árvores, 28 atributos candidatos por divisão e índice de Gini (Seção IV-A), e uma regressão logística que combina os três, com o \texttt{StackingClassifier} da biblioteca mlxtend (Seção IV-B). O artigo escolhe o Random Forest por ser um dos métodos supervisionados mais acurados (Seção IV-A) e justifica a profundidade 5 pela rapidez na predição (Seção IV-B); o \texttt{TreeExplainer} aparece na linha 8 do Algoritmo 1.

\textbf{Explicação.} Valores SHAP \cite{lundberg2017}: importância global, gráficos de dependência e explicações locais, em um painel interativo (Seção VI).

\textbf{Pontos em aberto.} O artigo dá duas indicações para a profundidade das árvores: profundidade máxima 5 (Seção IV-B) e \textit{variable tree depth} (linha 3 do Algoritmo 1). Também não informa a limpeza dos dados, a seed, os parâmetros do SMOTE, com que dados e com que entrada o meta-classificador é treinado, nem a média usada nas métricas. Cita uma \textit{one-sided selection} com o SMOTE sem descrevê-la (Seção III-B), e ela não foi aplicada: o Non-DoH só é dividido em três partes. Cita uma busca de hiperparâmetros com \texttt{GridSearchCV} em 10 folds sem listar a grade (Seção III-D); a reprodução usa os hiperparâmetros finais que o artigo informa, sem busca. O repositório público dos autores traz o painel e um script com um Random Forest único, sem o empilhamento; por isso o sistema foi reimplementado a partir do texto.

\section{Solução proposta pela equipe para melhorar a solução do artigo de referência}
\label{sec:equipe}

\textbf{O que se espera melhorar.} Em cada subconjunto, 92,09\% do Benign-DoH é sintético, e o empilhamento acrescenta um estágio em que a classe pode ser perdida (Seção \ref{sec:res-cira}). A hipótese, escrita antes da execução, é que um modelo mais simples, sem amostras sintéticas, tem F1 macro maior que o do empilhado, por menos de 1 ponto.

\textbf{Componentes.} A modificação, chamada M1M2, é um Random Forest único, treinado no treino inteiro, sem SMOTE e sem meta-classificador. O desbalanceamento é tratado com peso de classe: a classe $c$, com $n_c$ dos $n$ fluxos do treino e $K=3$ classes, recebe
\begin{equation}
w_c = \frac{n}{K\, n_c}.
\label{eq:peso}
\end{equation}
Os hiperparâmetros (número de árvores, profundidade máxima e atributos por divisão) são escolhidos em cada seed entre oito combinações, por validação cruzada de 5 folds em uma subamostra de 25\% do treino, pelo maior F1 macro; o teste não participa. M1 é o passo intermediário: só troca o SMOTE e o empilhamento pelo peso de classe e mantém os hiperparâmetros do artigo; M2 é a seleção de hiperparâmetros, e M1M2 reúne os dois. Nas tabelas e figuras, A é o sistema do artigo com profundidade variável, A-prof5 é o mesmo com profundidade 5, e M1-prof5 é M1 com profundidade 5. A entrada e a saída são as do sistema original.

\textbf{Robustez.} O artigo publica que \texttt{Duration} é o atributo mais importante e que acima de cerca de 40~s ele empurra a decisão para Malicious-DoH (Fig. 5 e 6a). Um atacante que leia isso pode encerrar e reabrir a conexão do túnel. Simulamos esse atacante dividindo \texttt{Duration}, \texttt{FlowBytesSent} e \texttt{FlowBytesReceived} dos fluxos maliciosos do teste por um fator $k$, sem retreinar o modelo. É uma perturbação no espaço de atributos: nenhum tráfego foi gerado.

\section{Metodologia}

\subsection{Dados usados pelo artigo de referência}

O CIRA-CIC-DoHBrw-2020 \cite{montazeri2020} é o conjunto do artigo e é público. Non-DoH é navegação HTTPS; Benign-DoH é a resolução DoH dos navegadores; Malicious-DoH são túneis de dns2tcp, DNSCat2 e iodine. O artigo não descreve a limpeza: remover os fluxos com valor ausente reproduz as contagens da Tabela I nas três classes (Tabela \ref{tab:cira}). Os cinco identificadores (endereços, portas e horário) não entram no modelo.

O split é estratificado, com 10\% para teste e seed 42. Não há conjunto de validação separado: a validação é cruzada, com 10 folds estratificados dentro do treino, como no artigo; a Tabela \ref{tab:cira} traz o tamanho de um fold. O teste tem 115\,911 fluxos, um a mais que o do artigo, por arredondamento da biblioteca.

\begin{table}[!t]
\centering
\caption{Fluxos por classe do CIRA-CIC-DoHBrw-2020. Validação é o menor dos 10 folds da validação cruzada sobre o treino. As contagens do artigo são as somas das linhas das Figs. 4a e 4b. Seed 42.}
\label{tab:cira}
\tabelacoluna{dados_cira_contagens}
\end{table}

Três características medidas limitam o que qualquer resultado neste conjunto permite concluir: (i) o Malicious-DoH foi capturado em 10 máquinas, de 18/03 a 01/04/2020, e as outras classes em outras 4, de 09/12/2019 a 14/01/2020, sem nenhuma máquina nem dia em comum; (ii) 13,67\% dos fluxos do teste (seed 42) têm vetor de atributos idêntico ao de um fluxo do treino, e 326 vetores ocorrem em mais de uma classe; (iii) 298\,766 fluxos têm o valor $-10$ em alguma coluna de assimetria, marcador que o extrator grava quando o desvio padrão é zero.

\subsection{Novo conjunto de dados escolhido}

O segundo conjunto é o combinado CIRA-CIC-DoHBrw-2020 com DoH-Tunnel-Traffic-HKD (HKD) \cite{mitsuhashi2023}, referência externa à bibliografia do artigo. Foi escolhido por três motivos: tem os mesmos atributos, extraídos com o DoHLyzer, com as mesmas colunas; traz túneis de três ferramentas ausentes do CIRA (dnstt, tcp-over-dns e tuns), capturados em outras máquinas, em 2021; e é o único com as três classes, pois o HKD sozinho só tem tráfego malicioso e não permite treinar o sistema.

Duas ressalvas: Non-DoH e Benign-DoH do combinado são os do CIRA, de modo que só a classe maliciosa ganha tráfego novo; e o arquivo publicado repete 20 vezes cada fluxo do HKD, o que põe cópias idênticas no treino e no teste. Por isso o conjunto principal é o combinado sem réplicas (Tabela \ref{tab:segundo}), com o mesmo protocolo de split e validação. Ao lado ficam o combinado como publicado e a transferência, em que o sistema treinado no CIRA é avaliado nos 5258 fluxos do HKD.

\begin{table}[!t]
\centering
\caption{Fluxos por classe do segundo conjunto de dados: combinado como publicado e sem réplicas. Validação é o menor dos 10 folds. O HKD isolado só serve de teste. Seed 42.}
\label{tab:segundo}
\tabelacoluna{dados_segundo_contagens}
\end{table}

\subsection{Métricas de avaliação}

Para a classe $c$, com verdadeiros positivos $TP_c$, falsos positivos $FP_c$, falsos negativos $FN_c$ e verdadeiros negativos $TN_c$:
\begin{equation}
P_c = \frac{TP_c}{TP_c + FP_c},\quad
R_c = \frac{TP_c}{TP_c + FN_c},\quad
F1_c = \frac{2 P_c R_c}{P_c + R_c},
\end{equation}
\begin{equation}
FPR_c = \frac{FP_c}{FP_c + TN_c},\qquad
F1_{\text{macro}} = \frac{1}{K}\sum_{c=1}^{K} F1_c .
\end{equation}
A acurácia é a fração de fluxos corretos. Com 1,70\% do teste em Benign-DoH, ela quase não registra a perda dessa classe; por isso as métricas por classe vêm antes das médias, e a média macro, que pesa as três classes por igual, é a principal. O recall de Malicious-DoH mede os túneis que passam sem alerta, e o FPR de Malicious-DoH contra o resto, os alarmes falsos. A AUC é a um-contra-o-resto com média macro. A distância à matriz do artigo é a soma das diferenças absolutas, célula a célula.

Experimentos: (1) reprodução nas duas leituras de profundidade (5 e variável), com seed 42; (2) os três modelos de comparação da Tabela II; (3) sensibilidade a cinco pontos em aberto; (4) protocolo corrigido, que troca a execução única da reprodução por 10 seeds, com comparação pareada no mesmo split e sem nenhuma seleção de modelo, e acrescenta o teste sem vetores repetidos e os folds por máquina (o SMOTE já era ajustado só no treino na reprodução); (5) explicabilidade com SHAP; (6) identificação da ferramenta de túnel; (7) os experimentos 1, 2 e 5 refeitos no segundo conjunto; (8) modificação e robustez, com 10 seeds. Nas comparações pareadas por seed, há diferença quando a média das diferenças passa, em módulo, do desvio padrão delas, regra fixada antes da execução. Nenhuma seed, hiperparâmetro ou regra de limpeza foi ajustada para aproximar um número do artigo.

Leituras adotadas onde o artigo é omisso: SMOTE com os parâmetros padrão da biblioteca; meta-classificador treinado no treino original, com os rótulos preditos pelos três bases; \texttt{TreeExplainer} aplicado a cada Random Forest base, porque ele não aceita o modelo empilhado. Os experimentos 5 a 8 usam como base a leitura de profundidade variável, porque um sistema que não prediz uma das classes não sustenta explicação, segundo conjunto nem modificação; a leitura de profundidade 5 é reportada ao lado.

\section{Resultados e discussões}

\subsection{Resultados da reprodução do artigo de referência}

\subsubsection{Nos dados usados pelo artigo de referência}
\label{sec:res-cira}

\begin{figure*}[!t]
\centering
\includegraphics[width=0.88\textwidth]{figures/matriz_confusao_teste_dupla.pdf}
\caption{Matrizes de confusão no teste: artigo (Fig. 4b) e reprodução nas duas leituras de profundidade. Em cada célula, o número de fluxos e, entre parênteses, a diferença para o artigo. Seed 42.}
\label{fig:matriz}
\end{figure*}

\begin{table}[!t]
\centering
\caption{Métricas no teste do CIRA, em \%: valores calculados da matriz da Fig. 4b do artigo, reprodução nas duas leituras e diferença em pontos percentuais. Seed 42; médias macro e ponderada nomeadas.}
\label{tab:reproducao}
\tabelacoluna{reproducao_metricas}
\end{table}

Todas as tabelas e figuras de resultado do artigo foram reproduzidas; este relatório mostra as principais, e as demais estão em \texttt{report/figures/} no repositório\footnote{\texttt{github.com/BrenoRev/Tecnicas\_Ataque\_Deteccao\_Intrusao}}: Fig. 2 (\texttt{fig2\us densidades\us dupla}), Fig. 4a (\texttt{matriz\us confusao\us validacao\us dupla}), Fig. 6 (\texttt{shap\us dependencia\us cira\us dupla}), Figs. 7 e 8 (\texttt{shap\us local\us cira\us dupla}) e Fig. 9 (\texttt{fig9\us ferramentas\us dupla}). Na validação cruzada, ao lado da Fig. 4a, a soma das diferenças absolutas é 43\,275 com profundidade 5 e 5147 com profundidade variável, e o F1 macro é 65,75\% e 96,57\%, contra 97,66\% calculado da matriz do artigo.

\textbf{Profundidade 5.} O sistema empilhado não prediz Benign-DoH em nenhum fluxo do teste (Fig. \ref{fig:matriz}, Tabela \ref{tab:reproducao}), e o mesmo ocorre nas dez seeds do protocolo corrigido. A acurácia de 97,79\% esconde a perda. A causa está no meta-classificador, não nos bases: sozinhos, os três Random Forests têm recall de Benign-DoH de 85,16\% a 85,97\%, com precisão de 27,89\% a 30,06\%. No treino, a combinação em que os três dizem Benign-DoH ocorre em 49\,247 fluxos, dos quais 31\,025 são Non-DoH; a regressão logística aprende a devolver Non-DoH para ela, e no teste isso atinge 5473 fluxos.

\textbf{Profundidade variável.} A matriz fica a 553 fluxos da Fig. 4b, contra 4711 na outra leitura; o F1 macro é 96,56\%, a 1,26 ponto do artigo, com recall de Benign-DoH maior e precisão menor. A proximidade não identifica a configuração dos autores: as duas leituras têm apoio no texto, e na análise de sensibilidade as cinco leituras alternativas dão somas de 469 a 1061, enquanto só a troca de seed move a soma de 509 a 653.

\begin{table*}[!t]
\centering
\caption{Metade superior da Tabela II do artigo ao lado da reprodução, em \%, com média macro em F1, precisão e recall (o artigo não diz que média usa) e AUC um-contra-o-resto macro. Seed 42.}
\label{tab:tabela2}
\tabelapagina{tabela2_superior_macro_dupla}
\end{table*}

\begin{table}[!t]
\centering
\caption{Metade inferior da Tabela II do artigo: resultados de Banadaki \cite{banadaki2020}, Jafar et al. \cite{jafar2021} e Ahakonye et al. \cite{ahakonye2022} no mesmo conjunto de dados, em \%. O artigo imprime sete das oito linhas em fração, aqui multiplicadas por 100, e a última já em percentual. O artigo declara que não são diretamente comparáveis.}
\label{tab:literatura}
\tabelacoluna{tabela2_literatura}
\end{table}

\textbf{Tabela II.} A Tabela \ref{tab:tabela2} mostra que nenhuma média reproduz a linha do modelo proposto: a Tabela II do artigo dá acurácia de 99,98\%, e a matriz que o próprio artigo publica na Fig. 4b dá 99,78\%. As duas fontes do artigo divergem entre si, e reportamos a distância às duas. Entre os modelos de comparação, a ordem do artigo se mantém em F1 macro (árvore de decisão, XGBoost, Random Forest), com o sistema de profundidade variável à frente e o de profundidade 5 atrás de todos. Os trabalhos da Tabela \ref{tab:literatura} reportam acurácia de até 99,98\% no mesmo conjunto de dados; como o artigo, não os reproduzimos, e eles não informam a média nem métricas por classe que permitam comparar o Benign-DoH.

\textbf{Protocolo corrigido.} Em dez seeds, o sistema de profundidade variável tem F1 macro de 96,56 $\pm$ 0,12\% e recall de Benign-DoH de 93,07 $\pm$ 0,52\%. Na média das dez seeds, 13,64\% do teste repete um vetor do treino; sem esses fluxos o F1 macro é 96,75 $\pm$ 0,14\% e os vereditos das comparações não mudam. Com folds por máquina de origem, o recall de Benign-DoH cai para 31,10 $\pm$ 9,42\% e o F1 macro para 77,89 $\pm$ 5,99\% (4 folds): a separação entre Non-DoH e Benign-DoH aprendida em três máquinas não vale na quarta. Esses folds não medem a generalização da detecção de túnel, porque nenhuma máquina gerou tráfego legítimo (Non-DoH e Benign-DoH) e malicioso.

\begin{figure}[!t]
\centering
\includegraphics[width=0.82\columnwidth]{figures/shap_importancia_cira.pdf}
\caption{Importância global dos 10 atributos de maior média do valor absoluto de SHAP para Malicious-DoH, equivalente à Fig. 5 do artigo, no Random Forest base 1, CIRA, nas duas leituras de profundidade. Amostra de 2000 fluxos por classe do treino; seed 42.}
\label{fig:shap}
\end{figure}

\textbf{Explicabilidade.} O artigo põe \texttt{Duration} em primeiro lugar na Fig. 5; aqui o primeiro é \texttt{PacketLengthMode}, com larga margem, e \texttt{Duration} é o segundo ou o terceiro (Fig. \ref{fig:shap}). A correlação de Spearman com o ranking do artigo é 0,62 (profundidade 5) e 0,51 (variável), com 5 e 7 atributos em comum entre os dez primeiros. A predominância de \texttt{PacketLengthMode} tem explicação nos dados: quatro valores do atributo cobrem 99,84\% do Malicious-DoH e ocorrem em 1 de 909\,555 fluxos legítimos. Uma regra com um só atributo teria esse recall, o que indica que o modelo pode separar as classes pela captura e não pelo túnel. Na dependência de \texttt{Duration}, a direção descrita pelo artigo se observa (acima de 40~s, têm SHAP positivo 97,37\% dos fluxos com profundidade variável e 99,53\% com profundidade 5); o sinal troca em 33,13~s na nossa amostra, de classes em partes iguais, e a diferença para os 40~s que o artigo lê na figura não é atribuível. Os valores SHAP explicam os bases, não a saída do empilhamento; o painel interativo foi implementado como material de demonstração.

\begin{table}[!t]
\centering
\caption{Identificação da ferramenta de túnel no CIRA, em \%, ao lado dos valores que a Seção VI-D do artigo chama de acurácia. O método é leitura da equipe; Dif. é o recall menos o valor do artigo. Seed 42; uma execução.}
\label{tab:ferramenta}
\tabelacoluna{ferramenta_metricas}
\end{table}

\textbf{Ferramenta de túnel.} O artigo dá três valores por ferramenta e não descreve o método (Seção VI-D). Treinamos o mesmo sistema só com os fluxos maliciosos e as três ferramentas como classes. Com profundidade variável, o recall fica a menos de 2 pontos dos valores do artigo; com profundidade 5, a até 20,28 (Tabela \ref{tab:ferramenta}). As 12 médias e desvios da legenda da Fig. 9 saem iguais nos seis algarismos impressos quando calculados com os arquivos por ferramenta sem a limpeza e com desvio populacional, e nenhuma com os dados limpos: é indício de que essa seção do artigo usou dados diferentes dos da Tabela I.

\subsubsection{No novo conjunto de dados escolhido}

\begin{table*}[!t]
\centering
\caption{Modelos de comparação da Tabela II e sistema proposto, nas duas leituras de profundidade, no teste do CIRA e do combinado sem réplicas, em \%. Seed 42; média macro.}
\label{tab:segundo-modelos}
\tabelapagina{segundo_baselines_dupla}
\end{table*}

\begin{table}[!t]
\centering
\caption{Recall de Malicious-DoH por ferramenta, em \%, no segundo conjunto: transferência (treino no CIRA, teste no HKD), retreino no combinado publicado e no sem réplicas. n é o número de fluxos no teste. Seed 42.}
\label{tab:recall-ferramenta}
\tabelacoluna{segundo_recall_ferramenta}
\end{table}

\textbf{Transferência.} O sistema treinado no CIRA detecta 1,81\% dos 5258 fluxos do HKD (1,71\% com profundidade 5): nenhum de dnstt e 6,20\% de tuns (Tabela \ref{tab:recall-ferramenta}). É o achado principal do trabalho. A regra de quatro valores de \texttt{PacketLengthMode}, que detecta 99,84\% no CIRA, detecta 0\% no HKD: o valor mais frequente lá é o de tráfego legítimo no CIRA. O resultado é compatível com um detector que aprendeu a assinatura das três ferramentas e da captura do CIRA, e não o comportamento de túnel. O que não foi medido é quanto da decisão do sistema nos fluxos do HKD vem desse atributo.

\textbf{Retreino.} No combinado sem réplicas, as métricas gerais repetem as do CIRA (Tabela \ref{tab:segundo-modelos}), o que é esperado, pois duas das três classes são as mesmas. Com profundidade 5, o sistema continua sem predizer Benign-DoH. O recall nas ferramentas do HKD é 99,42\% (510 de 513) com profundidade variável e 73,29\% com profundidade 5. Esse número mede as mesmas ferramentas em fluxos muito próximos dos do treino, não uma ferramenta nunca vista: em 505 dos 513, o fluxo do HKD mais próximo no treino está mais perto que qualquer fluxo Malicious-DoH do CIRA. No combinado publicado, os 10\,490 fluxos do HKD no teste têm cópia idêntica no treino, e o recall de 100\% mede memorização. Os três modelos de comparação têm F1 macro abaixo do sistema de profundidade variável também aqui.

\subsection{Resultados da proposta de melhoria do artigo de referência}

\begin{table*}[!t]
\centering
\caption{Modificação da equipe (M1M2: Random Forest único com peso de classe e seleção de hiperparâmetros) e passo intermediário (M1: sem a seleção; M1-prof5: M1 com profundidade 5) contra o sistema do artigo (A: profundidade variável; A-prof5: profundidade 5), no teste, em \%, com o tempo de treino em segundos. 10 seeds (0 a 9), média $\pm$ desvio padrão amostral; média macro.}
\label{tab:modificacao}
\tabelapagina{modificacao_metricas_dupla}
\end{table*}

\begin{figure}[!t]
\centering
\includegraphics[width=0.85\columnwidth]{figures/robustez_fragmentacao.pdf}
\caption{Recall de Malicious-DoH sob fragmentação simulada do fluxo, no CIRA: duração e bytes dos fluxos maliciosos do teste divididos por $k$, sem retreino. A e A-prof5: sistema do artigo com profundidade variável e com profundidade 5; M1M2: modificação da equipe. 10 seeds, média $\pm$ desvio padrão amostral.}
\label{fig:robustez}
\end{figure}

\textbf{Nos dados do artigo.} No CIRA, M1M2 tem F1 macro 0,50 ponto acima do sistema do artigo, nas dez seeds, e precisão de Benign-DoH 2,19 pontos maior (Tabela \ref{tab:modificacao}). A seleção escolheu a mesma combinação nas dez seeds, nos dois conjuntos: 100 árvores, sem limite de profundidade e raiz quadrada do número de atributos por divisão. O custo é o tempo de treino: 263,2~s contra 161,0~s do sistema do artigo (1,6 vez). O tempo do sistema do artigo no CIRA foi medido sob outra carga da máquina; reajustado nas mesmas seeds, é 112,5~s (2,3 vezes).

\textit{De onde vem o ganho.} M1, sem seleção de hiperparâmetros, já tem F1 macro 0,58 ponto acima do sistema do artigo, treina em 31,1~s e perde 0,50 ponto de recall de Benign-DoH; seu FPR de Malicious-DoH piora, de 0,0015\% para 0,0048\%. A diferença de F1 macro entre M1M2 e M1 é de $-$0,08 $\pm$ 0,15 ponto: os dois não se distinguem, e a seleção apenas troca precisão por recall de Benign-DoH. No protocolo corrigido, o modelo B, um Random Forest único com SMOTE, tem F1 macro de 95,71\%, abaixo dos 96,56\% do sistema do artigo: tirar só o empilhamento não deu ganho, e o ganho de M1 aparece quando o SMOTE sai e entra o peso de classe, com a ressalva de que o SMOTE de B age sobre o treino inteiro, e não sobre cada subconjunto. Com profundidade 5, o Random Forest único (M1-prof5) recupera o recall de Benign-DoH que o empilhamento perde (87,64\%), mas com precisão de 26,23\% e F1 macro de 78,66\%. No FPR de Malicious-DoH, M1M2 soma 2 falsos positivos nas dez seeds do CIRA, contra 14 do sistema do artigo, cerca de um fluxo por teste; essa melhora e a do recall de Benign-DoH não se repetem no combinado.

\textit{Robustez.} Com $k=2$, em que a maioria dos vetores perturbados continua coerente, o recall do sistema do artigo cai 3,19 pontos e o de M1M2, 0,34 (Fig. \ref{fig:robustez}). De $k=4$ em diante, o recall do sistema do artigo varia pouco entre seeds (desvio de até 0,19 ponto) e fica de 4,53 a 5,24 pontos abaixo do valor sem perturbação; o de M1M2 cai abaixo de 90\% em duas seeds ($k=4$) ou três ($k=8$ e $16$), e o desvio chega a 10,90 pontos. Nos fatores 8 e 16, mais de 96\% dos vetores têm tempo médio de pacote maior que a duração, relação que nenhum fluxo real tem. O resultado mostra que os modelos dependem da duração e dos bytes do fluxo, mas não ordena os dois em robustez nem estima a evasão de um atacante real.

\textbf{No novo conjunto de dados.} No combinado sem réplicas, M1M2 tem F1 macro 0,57 ponto acima do sistema do artigo e precisão de Benign-DoH 2,72 pontos maior, nas dez seeds, a 2,3 vezes o tempo de treino (265,1 contra 117,8~s). O combinado não é réplica independente do CIRA, pois só acrescenta os fluxos do HKD. Treinado no CIRA, M1M2 detecta de 2,13\% a 42,96\% dos fluxos do HKD conforme a seed. A variação impede qualquer conclusão sobre generalização, e não há par do sistema do artigo nas mesmas seeds.

\section{Conclusões e trabalhos futuros}

A reprodução depende de uma escolha que o artigo deixa em aberto. Com profundidade máxima 5, o sistema perde uma classe inteira no meta-classificador; com profundidade variável, chega perto da Fig. 4b, a 1,26 ponto de F1 macro. Os valores da Tabela II e do resumo do artigo não foram reproduzidos por nenhuma leitura, e divergem da matriz que o próprio artigo publica.

\textbf{Limitações do sistema do artigo.} (i) Não detectou as três ferramentas do HKD, de outra captura: recall de 1,81\%. (ii) Um único atributo, \texttt{PacketLengthMode}, separa quase todo o tráfego malicioso do CIRA, cuja captura se confunde com a classe. (iii) A separação entre Non-DoH e Benign-DoH não se sustenta entre máquinas. (iv) Um Random Forest único com peso de classe tem F1 macro maior que o empilhamento, que com profundidade 5 anula uma classe. (v) A explicação SHAP se refere aos bases, não à decisão final.

\textbf{Limitações deste trabalho.} A reprodução e o segundo conjunto têm uma execução, com seed 42; só o protocolo corrigido e a modificação têm dez. Os testes das seeds se sobrepõem, e por isso não usamos significância estatística. O segundo conjunto compartilha duas classes com o primeiro. A fragmentação é simulada no espaço de atributos. A modificação treina em 1,6 a 2,3 vezes o tempo do sistema do artigo e, sob fragmentação, seu recall cai abaixo de 90\% em duas ou três seeds. Ela melhora o F1 macro em meio ponto e não resolve nenhuma das limitações (i) a (iii).

\textbf{Trabalhos futuros.} Gerar tráfego de túnel fragmentado e reextrair os atributos; avaliar com tráfego legítimo e malicioso capturados nas mesmas máquinas; treinar com ferramentas do HKD e testar em uma ferramenta deixada de fora; explicar o modelo único, em que o \texttt{TreeExplainer} cobre a decisão inteira.

\begin{thebibliography}{10}

\bibitem{zebin2022}
T.~Zebin, S.~Rezvy, and Y.~Luo, ``An explainable AI-based intrusion detection system for DNS over HTTPS (DoH) attacks,'' \emph{IEEE Trans. Inf. Forensics Security}, vol.~17, pp. 2339--2349, 2022, doi: 10.1109/TIFS.2022.3183390.

\bibitem{lundberg2017}
S.~M. Lundberg and S.~Lee, ``A unified approach to interpreting model predictions,'' in \emph{Proc. 31st Int. Conf. Neural Information Processing Systems}, 2017, pp. 4768--4777.

\bibitem{aiello2013}
M.~Aiello, M.~Mongelli, and G.~Papaleo, ``Basic classifiers for DNS tunneling detection,'' in \emph{2013 IEEE Symp. Computers and Communications (ISCC)}, 2013, pp. 880--885.

\bibitem{preston2019}
R.~Preston, ``DNS tunneling detection with supervised learning,'' in \emph{2019 IEEE Int. Symp. Technologies for Homeland Security (HST)}, 2019, pp. 1--6.

\bibitem{montazeri2020}
M.~MontazeriShatoori, L.~Davidson, G.~Kaur, and A.~H. Lashkari, ``Detection of DoH tunnels using time-series classification of encrypted traffic,'' in \emph{2020 IEEE Intl Conf on Dependable, Autonomic and Secure Computing}, 2020, pp. 63--70.

\bibitem{banadaki2020}
Y.~M. Banadaki, ``Detecting malicious DNS over HTTPS traffic in domain name system using machine learning classifiers,'' \emph{J. Computer Sciences and Applications}, vol.~8, no.~2, pp. 46--55, 2020.

\bibitem{jafar2021}
M.~T. Jafar, M.~Al-Fawa'reh, Z.~Al-Hrahsheh \emph{et~al.}, ``Analysis and investigation of malicious DNS queries using CIRA-CIC-DoHBrw-2020 dataset,'' \emph{Manchester J. Artificial Intelligence and Applied Sciences}, vol.~2, pp. 65--70, 2021.

\bibitem{ahakonye2022}
L.~A.~C. Ahakonye, C.~I. Nwakanma, S.~O. Ajakwe \emph{et~al.}, ``Countering DNS vulnerability to attacks using ensemble learning,'' in \emph{2022 Int. Conf. Artificial Intelligence in Information and Communication (ICAIIC)}, 2022, pp. 7--10.

\bibitem{vekshin2020}
D.~Vekshin, K.~Hynek, and T.~Cejka, ``DoH Insight: Detecting DNS over HTTPS by machine learning,'' in \emph{Proc. 15th Int. Conf. Availability, Reliability and Security}, 2020, pp. 1--8.

\bibitem{behnke2021}
M.~Behnke, N.~Briner, D.~Cullen \emph{et~al.}, ``Feature engineering and machine learning model comparison for malicious activity detection in the DNS-over-HTTPS protocol,'' \emph{IEEE Access}, vol.~9, pp. 129\,902--129\,916, 2021.

\bibitem{ramakrishnan2022}
S.~Ramakrishnan and A.~Senthil~Rajan, ``Network attack detection with QNNBADT in minimal response times using minimized features,'' in \emph{Computer Networks and Inventive Communication Technologies}. Springer, 2022, pp. 563--579.

\bibitem{mitsuhashi2023}
R.~Mitsuhashi, Y.~Jin, K.~Iida, T.~Shinagawa, and Y.~Takai, ``Malicious DNS tunnel tool recognition using persistent DoH traffic analysis,'' \emph{IEEE Trans. Netw. Service Manag.}, vol.~20, no.~2, pp. 2086--2095, Jun. 2023, doi: 10.1109/TNSM.2022.3215681.

\end{thebibliography}

\end{document}
